\chapter{课程化优化与完整训练流程}

\section{模块定位}
第三章、第四章和第五章分别给出了 RCTA 的时序一致性、可靠性门控和双原型结构约束。它们解决的问题不同，所依赖的训练阶段也不同：基础对齐和源域分类可以从训练开始即发挥作用；伪标签和一致性约束需要模型具备初步判别能力；原型约束则更依赖稳定的类别结构。如果多个辅助目标在训练初期同时以较大权重注入，往往会破坏尚未形成的分类边界。

本章对应第一章提出的挑战四：多个辅助目标的作用时机不同，若同时强行注入可能破坏训练稳定性。为此，RCTA 将训练节奏本身纳入方法设计，一方面通过课程化权重调度控制各损失的进入时机，另一方面通过置信保护型模型选择和早停策略控制 checkpoint 选择过程，从而把“训练如何组织”也纳入完整方法的一部分。

\section{两级课程机制}
当前实现中的课程化策略并不是单一的“损失预热”，而是包含两个层次。

\subsection{门控课程}
第一层课程作用于第四章的样本筛选过程，即门控分数下界 $\tau(t)$ 与按类接纳比例 $\rho(t)$ 的渐进变化。训练早期使用更严格的分数门槛和更小的接纳比例，只允许少量高可靠样本进入伪监督；训练中后期逐步放宽门槛并增加接纳比例，使更多目标域样本参与训练。这个课程的主要作用对象是“\textbf{哪些样本可以被使用}”。

\subsection{损失课程}
第二层课程作用于各辅助损失的进入时机。若记任一辅助项的基础权重为 $\lambda$，开始步数为 $t_0$，预热长度为 $\Delta t$，则其课程化权重统一写为
\begin{equation}
\lambda(t)=
\begin{cases}
0, & t<t_0,\\
\lambda \cdot \min\left(\dfrac{t-t_0}{\Delta t},1\right), & t\ge t_0.
\end{cases}
\end{equation}
这个课程的主要作用对象是“\textbf{哪些目标以多大强度进入优化}”。

在当前 fixed-fold benchmark 配置中，伪标签损失从训练开始即预热，最大发电权重为 0.22，预热步数为 900；一致性损失从第 700 步开始，基础权重为 0.06，预热步数为 1200；原型损失从第 800 步开始，基础权重为 0.07，预热步数为 1200。这样的安排体现了“先建立可迁移底座，再引入目标域监督，最后强化类别结构”的训练逻辑。

\section{完整目标函数}
结合前述各章，RCTA 的完整优化目标可写为
\begin{align}
\mathcal{L}_{\mathrm{RCTA}}=
&\ \mathcal{L}_{cls}^{s}
+\lambda_{align}(t)\mathcal{L}_{align}
+\lambda_{mcc}\mathcal{L}_{mcc}
+\lambda_{pl}(t)\mathcal{L}_{pl}\nonumber\\
&+\lambda_{proto}(t)\mathcal{L}_{proto}
+\lambda_{cons}(t)\mathcal{L}_{cons}
+\lambda_{semi}\mathcal{L}_{semi-cons}
+\lambda_{u}\mathcal{L}_{u-ent}.
\end{align}
其中，$\mathcal{L}_{cls}^{s}$ 负责源域有标签监督，$\mathcal{L}_{align}$ 负责基础域对齐，$\mathcal{L}_{mcc}$ 用于抑制目标域类别混淆，$\mathcal{L}_{pl}$、$\mathcal{L}_{proto}$ 和 $\mathcal{L}_{cons}$ 分别由第四章、第五章和第三章给出，$\mathcal{L}_{semi-cons}$ 与 $\mathcal{L}_{u-ent}$ 则是第四章分层利用机制中的温和辅助项。

从章节逻辑上看，第三章提供目标域预测的平滑来源，第四章筛选并分层利用目标样本，第五章塑造类条件结构，本章则决定这些目标以怎样的节奏进入优化，并决定训练过程中如何选择更稳健的 checkpoint。换言之，本章并不是单独新增一个损失项，而是负责把前面三个模块真正组织成一个可训练、可收敛的整体。

\section{置信保护型模型选择与早停}
在无监督领域自适应中，训练目标本身只是问题的一半，checkpoint 选取同样会显著影响最终结论。若只根据目标域平均置信度选择模型，对抗式方法容易在尚未真正完成对齐时出现“置信度很高但预测结构已塌缩”的情况；若只根据源域评估准确率选择模型，又可能无法反映目标域适应程度。因此，当前 benchmark 为 RCTA 额外设计了置信保护型模型选择指标。

记源域评估准确率为 $\mathrm{Acc}_{s}^{eval}$，目标域平均置信度为 $\bar{q}_{t}$，域判别器准确率为 $\mathrm{Acc}_{d}$，则选择分数可写为
\begin{equation}
S_{\mathrm{guard}}=
w_s\mathrm{Acc}_{s}^{eval}
+w_q\bar{q}_t
-w_o\max(0,\bar{q}_t-q_{max})
-w_d\max(0,|\mathrm{Acc}_{d}-d^\star|-\delta_d),
\end{equation}
其中 $q_{max}$ 为允许的置信上限，$d^\star=0.5$ 表示理想的域混淆点，$\delta_d$ 表示可接受的域差容忍带。当前 benchmark 配置中，该指标使用 $w_s=0.80$、$w_q=0.20$、$w_o=1.00$、$w_d=0.25$，并设置信心上限为 0.90、域差容忍带为 0.18。

这个指标的直观含义是：在优先保证源域判别能力的前提下，鼓励目标域预测逐渐形成清晰结构，但惩罚过度自信和域对齐明显失衡的 checkpoint。与“训练完最后一个 epoch 直接取结果”相比，这种选择方式更能体现 RCTA 对稳定性的关注，也更符合 benchmark 比较对可复现性的要求。

\section{训练流程表}
算法~\ref{alg:rcta-training} 给出了 RCTA 的训练流程。与第三章到第五章分别介绍局部机制不同，本节的算法表强调的是各模块之间的先后关系和信息流动方向。

\begin{algorithm}[htbp]
\caption{RCTA 训练流程}
\label{alg:rcta-training}
\KwIn{源域批量 $\mathcal{B}_s$，目标域批量 $\mathcal{B}_t$，学生网络参数 $\theta_{\mathrm{stu}}$，教师网络参数 $\theta_{\mathrm{tea}}$，原型记忆 $C_s,C_t$}
\KwOut{更新后的学生网络、教师网络和双原型记忆}
初始化学生网络、教师网络和源/目标域原型记忆\;
\For{训练步 $t=1,2,\ldots,T$}{
  从源域和目标域采样 $\mathcal{B}_s,\mathcal{B}_t$\;
  计算源域分类损失 $\mathcal{L}_{cls}^{s}$ 和基础对齐损失 $\mathcal{L}_{align}$\;
  对目标域样本构造弱增强 $\tilde{x}_{w}^{t}$ 与强增强 $\tilde{x}_{s}^{t}$\;
  教师网络在弱增强视图上生成 $p^{tea}$ 和伪标签 $\hat{y}^{t}$\;
  学生网络在强增强视图上生成 $p^{stu}$，计算一致性信息\;
  根据置信度、逆熵、增强一致性和原型相似度计算可靠性得分 $r$\;
  按类别执行门控，得到门控集合 $\Omega_t$ 及可靠/半可靠/不可靠划分\;
  使用 $\Omega_t$ 计算伪标签损失，并对半可靠和不可靠样本施加温和正则\;
  结合参考原型计算源域吸引、目标域吸引、软权重目标吸引和原型分离损失\;
  根据课程函数得到 $\lambda_{pl}(t)$、$\lambda_{cons}(t)$、$\lambda_{proto}(t)$\;
  汇总 $\mathcal{L}_{\mathrm{RCTA}}$ 并更新学生网络参数 $\theta_{\mathrm{stu}}$\;
  通过 EMA 更新教师网络参数，并更新源/目标原型记忆\;
}
\end{algorithm}

\section{递进验证中的 D 模块含义}
在论文写作结构中，本章对应的 D 模块不仅仅是“又加了一个损失项”，而是“\textbf{把前三章模块组织成完整方法的训练机制}”。具体地，D 模块包含以下三层含义：

\begin{enumerate}
\item 将门控课程与损失课程统一起来，决定不同样本和不同损失在训练过程中何时进入优化；
\item 将模型选择与早停显式纳入方法定义，避免因 checkpoint 选择方式不透明而削弱结论可信度；
\item 通过算法流程表将第三章至第五章的机制串联起来，形成从样本生成、样本筛选到结构约束的闭环训练过程。
\end{enumerate}

正因为如此，实验章节中所谓的“A+B+C+D：完整 RCTA”并不是单纯在 A+B+C 上再叠加一个附加权重，而是采用本章给出的完整训练节奏、选择指标和更新流程来运行整个系统。

\section{本章小结}
本章围绕课程化优化与完整训练流程展开，说明了门控课程、损失课程、完整目标函数、置信保护型模型选择指标以及算法流程表。与前三章关注局部模块不同，本章的核心作用是把局部机制组织成一个可训练、可选择、可复现的完整方法，从而为第七章的 benchmark 运行与结果分析提供稳定执行基础。
